\documentclass[10pt,a4paper]{amsart}
\usepackage[margin=19mm]{geometry}
\usepackage{amsmath,amssymb,mathtools}
\usepackage[scr=rsfs,cal=euler]{mathalfa}
\usepackage{xfrac}
\usepackage[unicode,hidelinks,bookmarks=false]{hyperref}
\newcommand{\ie}{i.e.\ }
\newcommand{\cf}{cf.\ }
\newcommand{\resp}{respectively\ }
\newcommand{\Subsection}{\subsection}
\providecommand{\nobreakdash}{\nobreak\hbox{-}}
\setlength{\emergencystretch}{2em}
\multlinegap0pt
\usepackage{amssymb}
\usepackage{mathtools}
\usepackage[shortlabels]{enumitem}
\usepackage{bm}
\usepackage{tcolorbox}
\usepackage{xcolor}
\usepackage{float}
\usepackage{csquotes}
\newtheorem{lemma}{Lemma}
\newcommand\C{\mathbb{C}}
\DeclareMathOperator{\Hom}{Hom}
\newcommand{\OO}{\mathcal{O}}
\title{The logarithmic leaf complex and foliated d-semistability}
\author{Mauricio Corr\^ea, Pablo Perrella, Sebasti\'an Velazquez}
\date{Source: arXiv:2605.00285v1 (CC BY 4.0)}
\begin{document}
\maketitle

\noindent\emph{Source and licence.} The logarithmic leaf complex and foliated d-semistability. Version arXiv:2605.00285v1 (CC BY 4.0), distributed by arXiv under the Creative Commons Attribution 4.0 International licence, http://creativecommons.org/licenses/by/4.0/, as recorded in the arXiv licence metadata for this exact version and shown on the abstract page https://arxiv.org/abs/2605.00285v1. Full licence notice: \texttt{licenses/arxiv-2605.00285-CC-BY-4.0.txt}.\par\smallskip

\noindent\textit{Editorial scope.} Counted unit: the article's lemma lema:defLie (source0.tex lines 962--964). The article's complete original proof of it is delivered with it (source0.tex lines 965--1043, one \texttt{proof} environment of 79 lines). No mathematics is written, repaired or re-derived: the statement and the proof are the article's own lines, in the article's own notation, and no retained line is substituted or re-flowed.  No further source-local context is needed: every cross-reference of every retained line resolves inside the two delivered spans.  Nothing was omitted from any retained span, so the package carries no omission record.  No retained line contains a citation, so the wrapper carries no bibliography and no bibliography entry.  No retained line contains a figure environment, an image-inclusion command or drawing code, so no image is delivered and the package declares no asset dependency.  Editorial formatting only, in addition: an AMS \texttt{amsart} preamble that restates the article's own declarations and macros used by the retained lines, namely the environments \texttt{lemma} on separate counters, together with the standard packages those lines need.  The retained environments print in this document as Lemma 1 in order of appearance, whereas the article prints the same items with the numbering of its own full document.  The complete native source of the article as distributed in the arXiv e-print for this exact version is bundled unchanged as \texttt{sources/199548/source0.tex}.
\begin{lemma}
\label{lema:defLie}    Let $\mathcal H$ be a subalgebra of a Lie algebra $\mathcal G$ and $0\to J\to A'\to A\to 0$ be a small extension. If $\mathcal{H}_A\subseteq \mathcal G_A$ is a deformation of $\mathcal H\subseteq \mathcal G$ and $\mathcal G_{A'}$ is an extension of $\mathcal G_A$ over $A'$, then the obstruction to extend $\mathcal H_A$ to some $\mathcal H_{A'}\subseteq \mathcal G_{A'}$ lies in the group $H^2(\mathcal H,\mathcal G/\mathcal H)\otimes J$, where $H^2$ stands for the second Chevalley-Eilenberg cohomology of the $\mathcal H$-module $\mathcal{G}/\mathcal H$.
\end{lemma}

\begin{proof}
Let us first observe that since $\mathcal H$ is locally free,  so is $\mathcal H_{A}$. This means that we can locally extend the morphism 
of $\OO_{X_A}$-modules $f_A:\mathcal{H}_A\hookrightarrow \mathcal{G}_A$ to some $\OO_{X_{A'}}$-linear arrow $f_{A'}:\mathcal{H}_{A'}\hookrightarrow \mathcal{G}_{A'}$.
In the same manner, we can also extend the Lie bracket
$$ [\,,]_A: \bigwedge^2\mathcal H_{A}\to \mathcal H_{A}$$
to some $\OO_{X_{A'}}$-linear
$$ \{\,,\}_{A'}: \bigwedge^2\mathcal H_{A'}\to \mathcal H_{A'}.$$
If we were too lucky, these extensions would satisfy the equation 
$$ f_{A'}\{x,y\}=[f_{A'}(x),f_{A'}(y)]_{A'},$$
providing an extension of $\mathcal H_{A}\hookrightarrow \mathcal G_A$ to $A'$. 
If this was not the case, we could aim to modify them to force this equality. We will now see that the obstruction to do so lies in the claimed group.


Let us consider the map $b:\bigwedge^2\mathcal H_{A'}\to J\otimes_{A'}(\mathcal G_{A'}/\mathcal H_{A'})\simeq J\otimes_\C \mathcal G/\mathcal H$
defined as 
$$ b(x,y)=\overline{[f_{A'}(x),f_{A'}(y)]_{A'}-f_{A'}(\{x,y\}_{A'})}.$$
Since $m_{A'}\cdot J=0$, $b$ factors through some $\overline{b}:\bigwedge^2\mathcal H\to  J\otimes_\C \mathcal G/\mathcal H$:


there are natural isomorphisms 
$$\Hom_{A'}\bigg(\bigwedge^k \mathcal H_{A'}, J\otimes_{A'} \mathcal G_{A'}/\mathcal H_{A'}\bigg)\simeq \Hom_{\C}\bigg(\bigwedge^k\mathcal H,J\otimes_\C \mathcal G/\mathcal H\bigg). $$
If $v_1,v_2,v_3\in \mathcal H$ then
\begin{align*}
    \delta\overline{b}(v_1,v_2,v_3)&= \overline{\sum_{1\leq i\leq 3} (-1)^{i+1}[f(v_i),\overline{b}(\widehat{v}_i)]+\sum_{i<j} (-1)^{i+j}\overline{b}([v_i,v_j],v_k)}\\
    &=\overline{[f(v_1),\overline{b}(v_2,v_3)]}-\overline{[f(v_2),\overline{b}(v_1,v_3)]}+ \overline{[f(v_3),\overline{b}(v_1,v_2)]}\\ &\quad\quad- \overline{b}([v_1,v_2],v_3)+ \overline{b}([v_1,v_3],v_2) - \overline{b}([v_2,v_3],v_1).
\end{align*}
In order to compute this, we can pick liftings $v'_i\in \mathcal H_{A'}$ of $v_i\in \mathcal H=\mathcal H_{A'}/\mathfrak m_{A'}\mathcal H_{A'}$, so that
\begin{align*}
\big[f(v_i),\overline{b}(v_j,v_k)\big]&=\overline{\big[f_{A'}(v'_i),[f_{A'}(v'_j),f_{A'}(v_k)]_{A'}-f_{A'}(\{v'_j,v'_k\}_{A'})\big]}_{A'}.
\end{align*}
Putting the last two expressions together and using the fact that $[ \,,]_{A'}$ satisfies de Jacobi identity gives $\delta \overline{b}=0$.

If we had chosen a different lifting for the bracket, say $\{\,,\}'_{A'}=\{\,,\}_{A'}+\eta$ for some $\eta:\bigwedge^2\mathcal H_{A'}\to J\otimes_{A'}\mathcal H_{A'}$, then 
\begin{align*}
[f_{A'}(x),f_{A'}(y)]_{A'} - f_{A'}\{x,y\}
&=[f_{A'}(x),f_{A'}(y)]_{A'}-f_{A'}(\{x,y\}'_{A'}) \\
&=[f_{A'}(x),f_{A'}(y)]_{A'}-f_{A'}(\{x,y\}_{A'})-f_{A'}(\eta(x,y)) \\
&=b(x,y)-f_{A'}(\eta(x,y)),
\end{align*}
which indeed coincides with $b$ after passing to the quotient $\mathcal G_{A'}/\mathcal H_{A'}$.

On the other hand, a different choice of $f_{A'}$ changes $\overline{b}$ by a Chevalley--Eilenberg boundary: for a morphism of the form $f'_{A'}=f_{A'}+h$, with $h:\mathcal H_{A'}\to J\otimes\mathcal G_{A'}$, we get
\begin{align*}
b'(x,y)
&=
[f'_{A'}(x),f'_{A'}(y)]_{A'}-f'_{A'}(\{x,y\}_{A'}) \\
&=
[f_{A'}(x)+h(x),\,f_{A'}(y)+h(y)]_{A'}
   -f_{A'}(\{x,y\}_{A'})-h(\{x,y\}_{A'}) \\
&=b(x,y)+[f_{A'}(x),h(y)]_{A'}-[f_{A'}(y),h(x)]_{A'}
   -h(\{x,y\}_{A'}).
\end{align*}
On $J\otimes \mathcal G_{A'}/\mathcal H_{A'}\simeq J\otimes_\C(\mathcal G/\mathcal H)$, this becomes
\[
\bar b'(x,y)
=
\bar b(x,y)
+\overline{[f(x),h(y)]}
-\overline{[f(y),h(x)]}
-\overline{h([x,y])},
\]
because $\{x,y\}_{A'}$ reduces to $[x,y]\in\mathcal H$.
Observe that the right-hand side is
$\bar b'(x,y)=\bar b(x,y)+(\delta\bar h)(x,y)$,
where $\bar h:\mathcal H\to J\otimes_\C(\mathcal G/\mathcal H)$ denotes the class of $h$. This shows that the obstruction is a well-defined class in
$
H^2(\mathcal H,\mathcal G/\mathcal H)\otimes J.
$
Moreover, this obstruction is zero precisely when $\overline{b}=\delta \overline{h} $ for some $\overline{h}:\mathcal H\to J\otimes_\C(\mathcal G/\mathcal H)$. In this case, choosing an $\OO_{X_{A'}}$-lifting
$$
h:\mathcal H_{A'}\to J\otimes_{A'} \mathcal G_{A'},
$$
we see by the argument above that $f'_{A'}=f_{A'}-h$
is a $\OO_{X_{A'}}$-linear morphism that satisfies
\[
[ f'_{A'}(x),f'_{A'}(y)]_{A'} \in f'_{A'}(\mathcal H_{A'})
\]
for every $x,y$, hence defining an extension of $\mathcal H_{A'}\subseteq \mathcal G_{A'}$ of $\mathcal H_A\subseteq \mathcal G_A$.
\end{proof}

\end{document}
